\section{Additional Benchmark Details}
% APPENDIX: Detailed task descriptions, environment setup, scoring methodology

\subsection{Task Descriptions}

\textbf{Task~0: Environment Setup.}
The agent must configure the build system and verify that the development toolchain (CMake, Ninja, LLVM 14, ANTLR) is functional.
Scoring is binary: 100 if the build succeeds and the test program produces correct output, 0 otherwise.

\textbf{Task~1: Lexer Implementation.}
The agent implements a lexical analyzer (\texttt{lex.l}) that tokenizes preprocessed C source code.
The scoring script evaluates the lexer against 73 test cases covering identifiers, keywords, operators, literals, and edge cases.

\textbf{Task~2: Parser Implementation.}
The agent implements a syntax analyzer (\texttt{parser.y}) that constructs an AST/ASG from the token stream produced by Task~1.
Scoring evaluates parse correctness across 73 test inputs.

\textbf{Task~3: IR Generation.}
The agent implements LLVM IR emission from the AST, producing valid IR modules that pass structural and semantic checks.

\textbf{Task~4: IR Optimization.}
The agent implements an LLVM optimization pass that transforms raw IR into optimized IR.
This is the highest-weighted task ($w_4 = 0.30$) reflecting its central role in compiler performance.

\textbf{Task~5: Backend Code Generation.}
The agent implements RISC-V 64-bit assembly emission from optimized IR.
Scoring uses the ``classic'' evaluation mode (\texttt{task5-classic-score}) which tests assembly correctness on a comprehensive benchmark suite.

\subsection{Environment Configuration}

Each evaluation workspace includes:
\begin{itemize}[nosep]
\item Ubuntu 22.04 base environment
\item CMake 3.20+ and Ninja build system
\item LLVM 14 development libraries
\item ANTLR 4 runtime
\item pybind11 for Python--C++ bindings
\item Python 3.10+ for scoring scripts
\end{itemize}

Workspaces are created as independent copies of the \yatcc{} repository to prevent cross-run interference.
CMake configuration is performed once per workspace, and build artifacts persist across tasks within a single run.

\subsection{Scoring Methodology}

Each task's scoring script (\texttt{score.py} or \texttt{score-classic.py}) performs:
\begin{enumerate}[nosep]
\item Build the target: \texttt{cmake --build build -t task\textit{k}}
\item Run the scorer: \texttt{cmake --build build -t task\textit{k}-score}
\item Parse the output: \texttt{score.json} or \texttt{score.txt}
\item Extract the total score and per-test details
\end{enumerate}

Scores are normalized to 0--100, with the passing threshold set at 60\%.